\subsection{Independent cross-program audit and the continuous-uncertainty gap}
\label{sec:p11-audit}
An independent reproducibility workflow cross-checked pinned feeder source hashes and archived numerical results, replayed the 1,200-case feeder-proxy study and 1,440-case joint frequency--relay experiment, and checked the reduced-order frequency envelope against separate fourth-order Runge--Kutta integration of constant-input trajectories. The archived Prompt 10 study has 18 finite-corner approvals across three relay reset semantics; the physical-evidence gate authorizes zero switching actions. These outcomes must not be conflated.

An additional audit tested 1,458 interior combinations of independently derated primary and backup current and breaker delay, drawn from the corner-approved configurations. No interior failure was observed. This is a finite falsification exercise, \emph{not} a proof that extrema occur at the corners of the entire continuous uncertainty box for nonlinear inverse-time relay dynamics with pickup and reset. The finite-corner screen is therefore not presented as a universal robust protection certificate outside explicitly established monotonic conditions. In the constant-current setting, monotonicity of the inverse-time law supplies the relevant endpoint reasoning; in interrupted-current settings, relay memory, dropout and current-direction assumptions must be examined separately.

The cross-program audit also identifies a historical validation weakness: an earlier inverse-time synthetic script increments its purported robust-unsafe counter by zero without evaluating an independent unsafe predicate. Its reported zero is not independent evidence of soundness. Later independent-corner and time-domain checks address narrower, explicitly stipulated models and must not be used to retroactively validate that counter.

The audit establishes consistency of archived results and code under pinned assumptions. It does not independently reconstruct every published figure from source data, verify the existence and metadata of all cited publications against publisher records, establish phase-specific energized paths, reproduce electromagnetic transients, or validate commercial relay firmware. No installed device settings, breaker timing, inverter current limiter or physical restoration action has been independently validated. The conditional logical conjunction of safety requirements is not asserted to be a novel protection principle by itself.
